\begin{abstract}
Ontologies pair a class hierarchy with names that embeddings can compare, and a
natural question combines the two: which terms are closest in meaning to a given
term \emph{and lie under a given class}?
Vector databases answer it by passing the index a filter, and approximate indexes
silently lose recall when the filter is selective: on four ontologies, FAISS
indexes asked for the ten nearest terms under a class return on average 61\% (IVF)
and 64\% (HNSW) of the right answers, and 2\% (HNSW) for a query that lies outside a
small class; a 16 times larger beam raises that to only 20\%.
This paper presents a vector search that lives inside an ontology store and uses
the hierarchy as the filter.
(1) We prove that no linear order of points keeps nearest neighbors close, and
measure how few a sorted index of an ontology puts within a short scan.
(2) We label the class hierarchy with postorder intervals in the store's own
B+ tree, so that a class's members are a few range scans and membership is a
binary search, and we store IVF posting lists as key ranges of the same tree.
(3) We show that in an ontology the class filter is correlated with the vectors, so
the independence assumption behind filtered-search costing fails, and give
four filtering strategies, a cost model for the crossover between resolving the
class and testing it per row, and a rule that picks between them.
(4) We evaluate against FAISS, hnswlib, and DuckDB on the Gene Ontology, the Human
Phenotype Ontology, Uberon, and ChEBI (218,723 terms), with queries that are
stored terms and queries that are not.
Across 958 queries our rule returns recall 0.999 and never falls below 0.99, with a
median of 26\,$\mu$s; for classes of at most ten members it is 33 times faster than
an exact FAISS scan with an ID selector (2.1 against 70\,$\mu$s).
\end{abstract}
